\begin{tikzpicture}[font=\figurefont\small,
  box/.style={draw=BookRule,fill=BookPaper,align=center,
    text width=24mm,minimum height=12mm},
  dot/.style={circle,fill=FigureModel!18,draw=BookRule,minimum size=5mm},
  arr/.style={-{Stealth[length=2mm]},draw=BookInk,line width=.8pt},
  train/.style={arr,dashed,draw=BookTeal}]
  \node[box,fill=FigureInput!12] (root) at (0,0) {真实起点\\$\bm z_0=e(s_t)$};
  \node[dot] (a1) at (3,1.3) {};
  \node[dot] (a2) at (5,1.3) {};
  \node[dot] (b1) at (3,0) {};
  \node[dot] (b2) at (5,0) {};
  \node[dot] (c1) at (3,-1.3) {};
  \node[dot] (c2) at (5,-1.3) {};
  \node[box,fill=FigureOutput!12] (q) at (8,0) {累计奖励\\加末端$Q$};
  \node[box,fill=FigureOutput!12] (act) at (12,0) {选最高分序列\\执行首动作};
  \foreach \x/\y in {a1/a2,b1/b2,c1/c2}{
    \draw[arr] (root.east) -- (\x);
    \draw[arr] (\x) -- (\y);
    \draw[arr] (\y) -- (q.west);
  }
  \draw[arr] (q) -- (act);
  \node[box,fill=FigureModel!12] (prior) at (4,3) {策略先验$\pi$\\提出部分候选};
  \draw[arr] (prior.south) -- (4,1.8) -- (a2.north);
  \node at (8.4,2.3) {其余候选由轨迹扰动生成};
  \node[box,fill=FigureInput!12] (data) at (0,-3.4) {真实短序列\\$(s,a,r,s',c)$};
  \node[draw=BookRule,fill=FigureLoss!12,align=center,
    text width=73mm,minimum height=13mm] (loss) at (7,-3.4)
    {潜在一致性$+$奖励预测$+$TD目标\\先验另通过提高评论家评分学习};
  \draw[train] (data) -- (loss);
  \draw[train] (loss.north) -- (7,-2.1) -- (4,-2.1) -- (c2.south);
  \node[align=center,text width=126mm] at (6,-5)
    {实线：本轮规划与执行；虚线：真实数据训练。\\策略先验缩小搜索负担，但不是唯一执行器。};
\end{tikzpicture}
